\documentclass[12pt]{article}
\usepackage{amsmath, amssymb, graphicx, hyperref, geometry, listings, xcolor}
\geometry{a4paper, margin=1in}

\definecolor{codegreen}{rgb}{0,0.6,0}
\definecolor{codegray}{rgb}{0.5,0.5,0.5}
\definecolor{codepurple}{rgb}{0.58,0,0.82}
\definecolor{backcolour}{rgb}{0.95,0.95,0.92}

\lstdefinestyle{mystyle}{
    backgroundcolor=\color{backcolour},   
    commentstyle=\color{codegreen},
    keywordstyle=\color{magenta},
    numberstyle=\color{codegray},
    stringstyle=\color{codepurple},
    basicstyle=\ttfamily\footnotesize,
    breakatwhitespace=false,         
    breaklines=true,                 
    captionpos=b,                    
    keepspaces=true,                 
    numbers=left,                    
    numbersep=5pt,                  
    showspaces=false,                
    showstringspaces=false,
    showtabs=false,                  
    tabsize=2
}
\lstset{style=mystyle}

\title{\textbf{Transformers from Scratch: Part 2 - Pure PyTorch Implementation}}
\author{Build in Public Series}
\date{\today}

\begin{document}
\maketitle

\begin{center}
    \includegraphics[width=\textwidth]{Post-2-Banner.png}
\end{center}

\section{Introduction}
Welcome to Part 2 of our series on Transformers. In Part 1, we outlined the mathematical foundations of Self-Attention and Positional Encodings. Now, it is time to turn those equations into running code.

Rather than relying on the high-level \texttt{nn.Transformer} module, we are building the entire sequence-to-sequence model from scratch in PyTorch. This requires precise management of tensor dimensions, multi-head reshapes, and causality masks.

\section{Implementing Scaled Dot-Product Attention}
The attention mechanism requires projecting our inputs, performing batched matrix multiplications, and applying masking.

\begin{lstlisting}[language=Python]
class ScaledDotProductAttention(nn.Module):
    def forward(self, q, k, v, mask=None):
        d_k = q.size(-1)
        scores = torch.matmul(q, k.transpose(-2, -1)) / math.sqrt(d_k)
        
        if mask is not None:
            scores = scores.masked_fill(mask == 0, -1e9)
            
        attention_weights = F.softmax(scores, dim=-1)
        output = torch.matmul(attention_weights, v)
        return output, attention_weights
\end{lstlisting}

Notice the use of \texttt{-1e9} when applying the mask. When we pass this tensor through the softmax activation, those values accurately converge to zero, ensuring the model cannot attend to padded or future tokens.

\section{Constructing the Encoder and Decoder Layers}
The Transformer architecture is symmetric, built of layered building blocks. 
The \textbf{Encoder Layer} consists of:
\begin{enumerate}
    \item Multi-Head Self-Attention
    \item Residual Connection \& Layer Normalization
    \item Position-wise Feed-Forward Network
    \item Residual Connection \& Layer Normalization
\end{enumerate}

The \textbf{Decoder Layer} is slightly more complex. It adds a \textit{Cross-Attention} block in the middle, which acts as the bridge between the encoder and the decoder. In cross-attention, the Queries come from the previous Decoder layer, while the Keys and Values are provided by the Encoder's final output.

\section{The Future-Blind Mask}
One of the trickiest parts of the implementation is the Look-Ahead Mask (or causal mask) in the decoder. When generating text sequentially, the model must be blind to future tokens during training (teacher forcing). We achieve this by applying a lower-triangular boolean mask matrix to the self-attention scores in the decoder.

\begin{lstlisting}[language=Python]
tgt_lookahead_mask = torch.tril(torch.ones((tgt_len, tgt_len))).bool()
\end{lstlisting}

\section{Up Next}
With the code complete, we now have a working Transformer. But how do we know it learns? In Part 3, we will set up a sequence-to-sequence training loop, visualize the loss curve, and extract the internal cross-attention tensors to visualize how the model routes information.
\end{document}
